\documentclass[10pt,a4paper]{amsart}
\usepackage{amsmath,amssymb,amsthm,graphicx}
\usepackage{verbatim}
\newenvironment{reference}{\comment}{\endcomment}
%\newenvironment{reference}{}{}
\newenvironment{slogan}{\comment}{\endcomment}
\newenvironment{history}{\comment}{\endcomment}
% For commutative diagrams we use Xy-pic
\usepackage[all]{xy}

% We use 2cell for 2-commutative diagrams.
\xyoption{2cell}
% Theorem environments.
%
\theoremstyle{plain}
\newtheorem{theorem}[subsection]{Theorem}
\newtheorem{proposition}[subsection]{Proposition}
\newtheorem{lemma}[subsection]{Lemma}

\theoremstyle{definition}
\newtheorem{definition}[subsection]{Definition}
\newtheorem{example}[subsection]{Example}
\newtheorem{exercise}[subsection]{Exercise}
\newtheorem{situation}[subsection]{Situation}

\theoremstyle{remark}
\newtheorem{remark}[subsection]{Remark}
\newtheorem{remarks}[subsection]{Remarks}

\numberwithin{equation}{subsection}

% Macros
%
\def\lim{\mathop{\mathrm{lim}}\nolimits}
\def\colim{\mathop{\mathrm{colim}}\nolimits}
\def\Spec{\mathop{\mathrm{Spec}}}
\def\Hom{\mathop{\mathrm{Hom}}\nolimits}
\def\Ext{\mathop{\mathrm{Ext}}\nolimits}
\def\SheafHom{\mathop{\mathcal{H}\!\mathit{om}}\nolimits}
\def\SheafExt{\mathop{\mathcal{E}\!\mathit{xt}}\nolimits}
\def\Sch{\mathit{Sch}}
\def\Mor{\mathop{\mathrm{Mor}}\nolimits}
\def\Ob{\mathop{\mathrm{Ob}}\nolimits}
\def\Sh{\mathop{\mathit{Sh}}\nolimits}
\def\NL{\mathop{N\!L}\nolimits}
\def\CH{\mathop{\mathrm{CH}}\nolimits}
\def\proetale{{pro\text{-}\acute{e}tale}}
\def\etale{{\acute{e}tale}}
\def\QCoh{\mathit{QCoh}}
\def\Ker{\mathop{\mathrm{Ker}}}
\def\Im{\mathop{\mathrm{Im}}}
\def\Coker{\mathop{\mathrm{Coker}}}
\def\Coim{\mathop{\mathrm{Coim}}}

% Boxtimes
%
\DeclareMathSymbol{\boxtimes}{\mathbin}{AMSa}{"02}

%
% Macros for moduli stacks/spaces
%
\def\QCohstack{\mathcal{QC}\!\mathit{oh}}
\def\Cohstack{\mathcal{C}\!\mathit{oh}}
\def\Spacesstack{\mathcal{S}\!\mathit{paces}}
\def\Quotfunctor{\mathrm{Quot}}
\def\Hilbfunctor{\mathrm{Hilb}}
\def\Curvesstack{\mathcal{C}\!\mathit{urves}}
\def\Polarizedstack{\mathcal{P}\!\mathit{olarized}}
\def\Complexesstack{\mathcal{C}\!\mathit{omplexes}}
% \Pic is the operator that assigns to X its picard group, usage \Pic(X)
% \Picardstack_{X/B} denotes the Picard stack of X over B
% \Picardfunctor_{X/B} denotes the Picard functor of X over B
\def\Pic{\mathop{\mathrm{Pic}}\nolimits}
\def\Picardstack{\mathcal{P}\!\mathit{ic}}
\def\Picardfunctor{\mathrm{Pic}}
\def\Deformationcategory{\mathcal{D}\!\mathit{ef}}
\usepackage{cleveref}
\crefname{originallemma}{original lemma}{original lemmas}
\Crefname{originallemma}{Original lemma}{Original lemmas}
\makeatletter
\newlabel{lemma-base-change-locally-quasi-finite}{{lemma-base-change-locally-quasi-finite}{source}{}{sourcepointer-1}{}}
\newlabel{lemma-base-change-locally-quasi-finite@cref}{{[originallemma][0][]lemma-base-change-locally-quasi-finite}{[1][0][]source}{}{}{}}
\makeatother
\crefname{originallemma}{original lemma}{original lemmas}
\Crefname{originallemma}{Original lemma}{Original lemmas}
\makeatletter
\newlabel{lemma-composition-locally-quasi-finite}{{lemma-composition-locally-quasi-finite}{source}{}{sourcepointer-2}{}}
\newlabel{lemma-composition-locally-quasi-finite@cref}{{[originallemma][0][]lemma-composition-locally-quasi-finite}{[1][0][]source}{}{}{}}
\makeatother
\crefname{originallemma}{original lemma}{original lemmas}
\Crefname{originallemma}{Original lemma}{Original lemmas}
\makeatletter
\newlabel{lemma-semi-diagonal}{{lemma-semi-diagonal}{source}{}{sourcepointer-3}{}}
\newlabel{lemma-semi-diagonal@cref}{{[originallemma][0][]lemma-semi-diagonal}{[1][0][]source}{}{}{}}
\makeatother
\title{A locally quasi-finite morphism over a quasi-DM morphism is locally quasi-finite}
\author{The Stacks Project Authors}
\date{stacks-project:stacks-morphisms.tex:local-snapshot-20261012; GNU FDL 1.2 or later}
\begin{document}
\maketitle
\noindent\textit{Editorial note.} All delivered mathematical text is the author's own, copied from the pinned Stacks Project chapter source sources/stacks-algebraic-stacks-morphisms/pinned-original-source.tex; no statement or proof step was rewritten, reordered or removed.\par\medskip
\noindent\textit{Editorial note.} Preamble provenance: the Stacks Project keeps one single shared preamble for every chapter; this card takes it directly from the staged sources/stacks-algebraic-stacks-morphisms/preamble.tex as three verbatim spans (lines 8-12, 14-18 and 154-217), so no derived file was written and nothing was added or redefined. The wrapper class and the base packages are supplied by the assembler.\par\medskip
\noindent\textit{Editorial note.} Bibliography provenance: neither the counted statement nor its proof contains a citation (this lane rescans the whole span before writing the card), so the card carries no bibliography and no entry was reconstructed.\par\medskip
\noindent\textit{Editorial note.} Selection provenance: this card was opened under the file-based control contract (CONTROL.md), with the mechanical pre-card gate run on its statement and proof interval before the card was written.\par\medskip
\noindent\textit{Allowed standard prerequisites (editorial).}\begin{itemize}
\item The base change permanence of being locally quasi-finite. Invoked inside the counted proof (source lines 5229-5231). Makes the projection of the fibre product locally quasi-finite.
\item The statement that the semi-diagonal of a quasi-DM morphism is locally quasi-finite. Invoked inside the counted proof (source lines 5232-5233). Makes the factorization arrow locally quasi-finite.
\item The composition permanence of being locally quasi-finite. Invoked at the end of the counted proof (source lines 5234-5236). Combines the two factors into the assertion.
\end{itemize}\medskip
\ifdefined\hypertarget\hypertarget{sourcepointer-1}{}\fi
\noindent\textit{Source reference.} Original lemma lemma-base-change-locally-quasi-finite: sources/stacks-algebraic-stacks-morphisms/pinned-original-source.tex, source lines 5050--5050. Navigation pointer only: the counted text refers to the Stacks item lemma-base-change-locally-quasi-finite, which belongs to another chapter of the project and is not part of this card. The pointer carries the target's own label as its number because this chapter's staged snapshot does not contain the other chapter's tag table; its content is not reproduced here.\par\medskip
\ifdefined\hypertarget\hypertarget{sourcepointer-2}{}\fi
\noindent\textit{Source reference.} Original lemma lemma-composition-locally-quasi-finite: sources/stacks-algebraic-stacks-morphisms/pinned-original-source.tex, source lines 5160--5160. Navigation pointer only: the counted text refers to the Stacks item lemma-composition-locally-quasi-finite, which belongs to another chapter of the project and is not part of this card. The pointer carries the target's own label as its number because this chapter's staged snapshot does not contain the other chapter's tag table; its content is not reproduced here.\par\medskip
\ifdefined\hypertarget\hypertarget{sourcepointer-3}{}\fi
\noindent\textit{Source reference.} Original lemma lemma-semi-diagonal: sources/stacks-algebraic-stacks-morphisms/pinned-original-source.tex, source lines 845--845. Navigation pointer only: the counted text refers to the Stacks item lemma-semi-diagonal, which belongs to another chapter of the project and is not part of this card. The pointer carries the target's own label as its number because this chapter's staged snapshot does not contain the other chapter's tag table; its content is not reproduced here.\par\medskip
\begin{lemma}
\label{lemma-locally-quasi-finite-permanence}
Let $\mathcal{X} \to \mathcal{Y} \to \mathcal{Z}$ be morphisms
of algebraic stacks. Assume that $\mathcal{X} \to \mathcal{Z}$
is locally quasi-finite and $\mathcal{Y} \to \mathcal{Z}$ is quasi-DM.
Then $\mathcal{X} \to \mathcal{Y}$ is locally quasi-finite.
\end{lemma}

\begin{proof}
Write $\mathcal{X} \to \mathcal{Y}$ as the composition
$$
\mathcal{X} \longrightarrow
\mathcal{X} \times_\mathcal{Z} \mathcal{Y} \longrightarrow
\mathcal{Y}
$$
The second arrow is locally quasi-finite as a base change of
$\mathcal{X} \to \mathcal{Z}$, see
Lemma \ref{lemma-base-change-locally-quasi-finite}.
The first arrow is locally quasi-finite by
Lemma \ref{lemma-semi-diagonal}
as $\mathcal{Y} \to \mathcal{Z}$ is quasi-DM.
Hence $\mathcal{X} \to \mathcal{Y}$ is locally quasi-finite by
Lemma \ref{lemma-composition-locally-quasi-finite}.
\end{proof}

\end{document}
